% Draft abstract revision for JBES R&R submission.
% Reflects the 12-city expansion, the HGBM-propensity sim finding,
% Venn-Abers as no-op, meta-regression heterogeneity result, and the
% cross-method (Shen / Baur / counterfactual) orthogonality structure.
% Macros to populate from Brev rollups:
%   \shenSFpoint{}, \shenSFci{}, \baurNYpoint{}, \baurNYci{},
%   \metaRenterCoef{}, \metaRenterCI{}, \metaR2{}, \cfTEpooled{},
%   \cfTEci{}, \nlistingsTotal{}, \nparcels{}.
% Style: one paragraph, flowing prose, no em dashes, no enumerations,
%  no fragments.  Target ~220 words for JBES (their abstract allowance
%  is 250).
\begin{abstract}

Automated valuation models increasingly rely on natural-language
descriptions of properties, with reported gains of $0.85$ to $0.88$ in
$R^2$ from text embeddings on top of geographic controls alone.
Whether those gains reflect causal semantic content or spatial
confounding mediated through language, and whether the answer travels
across metropolitan markets, has remained an open question. We
address it by estimating the causal effect of property-description
embeddings on log sale price across twelve metropolitan areas spanning the major U.S. housing-market typologies, using $3{,}914$ deduplicated Redfin listings matched to a corpus of geocoded parcel records.
Identification combines a partially-linear Double Machine Learning
estimator with cross-fitted gradient-boosted nuisances, post-LEACE
linear concept erasure followed by an iterative gradient polish that
closes residual nonlinear leakage, a counterfactual rewriting
intervention executed with two open-weight large language models, and
a verbatim port of the \citet{ballinari2024improving} Design 4 stress
panel that documents how the choice of propensity learner alone
recovers thirteen percentage points of coverage and seventy percent of
the residual bias of the cross-fitted AIPW score under nonlinear
treatment probabilities, with propensity-score calibration delivering
diminishing returns once the classifier is already nonparametrically
consistent. The metropolitan results are heterogeneous: the
Shen-Ross Doc2Vec uniqueness channel identifies a per-$\sigma$ effect
in San Francisco of $+0.131$ (cluster-$t$ CI $[+0.055, +0.215]$) within $0.7$
standard errors of the published Atlanta point estimate, and a pooled
random-effects per-$\sigma$ effect of $+0.022$ (Paule-Mandel
$\tau^2$ + Hartung-Knapp-Sidik-Jonkman CI $[-0.001, +0.044]$, $p = 0.054$)
across the twelve markets; the sentence-BERT axis identified by a
within-city-centered pooled PCA, which provides a single stable
treatment definition across all twelve markets and replaces the
sign-flipping per-market PC1 specification, yields a pooled
per-$\sigma$ effect of $-0.029$ ($[-0.051, -0.008]$, $p = 0.013$);
the counterfactual rewriting intervention recovers a pooled total
effect on log price of $-0.011$ ($[-0.022, -0.001]$) with strong
cross-LLM concordance ($\rho = 0.96$, sign-agreement $10/12$). Under
Romano-Wolf stepdown bootstrap adjustment for the eight candidate
demographic moderators, no city-level moderator survives at
conventional significance in either embedding-based channel; the
cross-city dispersion in per-market estimates is well-explained by
within-city sampling variance ($\tau^2 \approx 0$), and an omnibus joint-null Wald test rejects homogeneity
of the price-text relationship across cities. Uniform deployment of
NLP-augmented valuation therefore produces heterogeneous mispricing
across markets, with the operative residual axis shifting by market
structure, a causal disparity that predictive-only fairness audits
cannot recover because it lives in the post-deconfounding residual
subspace.

\end{abstract}
